\uuid{1f22}
\exo7id{3654}
\titre{exo7 3654}
\auteur{quercia}
\organisation{exo7}
\datecreate{2010-03-10}
\isIndication{false}
\isCorrection{false}
\chapitre{Espace vectoriel}
\sousChapitre{Somme directe}
\cpge{mpsi-mp2i.matrices-applications-lineaires.a, pcsi.matrices-applications-lineaires.a, ptsi.matrices-applications-lineaires.a, mpsi-mp2i.espaces-vectoriels.d, pcsi.espaces-vectoriels.d, ptsi.espaces-vectoriels.d, mp-mpi.reduction-endomorphismes.a, pc.complements-espaces-vectoriels.a, psi.complements-espaces-vectoriels.a, pt.complements-algebre-lineaire.a, bcpst2.applications-lineaires-matrices.c}
\module{Algèbre}
\niveau{L1}
\difficulte{}

\contenu{
\texte{
Soit $E$ un $ K$-ev de dimension finie $n$ et
${\cal B} = ({\vec e}_1,\dots,{\vec e}_n)$ une base de $E$.

On note $F_i = \{ u \in \mathcal{L}(E)$ tq $\Im u \subset \text{vect}({\vec e}_i)\}$.
}
\begin{enumerate}
    \item \question{Caractériser matriciellement les éléments de $F_i$.}
    \item \question{Montrer que $F_1 \oplus F_2 \oplus \dots \oplus F_n = \mathcal{L}(E)$.}
\end{enumerate}
}
